\documentclass{article}
\usepackage[review]{log_2026}
\usepackage{booktabs}
\usepackage{multirow}
\usepackage{amsfonts}
\usepackage{graphicx}
\usepackage[numbers,compress,sort]{natbib}
\usepackage{xcolor}

% Placeholders for results still landing; every \pending must be resolved
% before the PDF builds for submission (check_submission greps for it).
\newcommand{\pending}[1]{\textcolor{red}{[PENDING: #1]}}

\title{Anchor, Don't Name: What Ontology Guidance Actually Buys in LLM-Built Knowledge Graphs}

\author{Anonymous Authors}

\begin{document}
\maketitle

\begin{abstract}
Practitioners hand LLM extractors an ontology expecting better knowledge
graphs. We measure what the ontology actually buys in a pre-registered study
on financial filings (FinDER): five extraction schemas from schema-free to
full FIBO compared head-to-head, with the two decisive conditions replicated
at scale across three extractor LLMs and three stratified samples.
(1)~Ontology guidance does not improve extraction: it \emph{reduces}
cross-model agreement on entity names and does not raise coverage of
answer-bearing facts; its measurable contribution is \emph{detectability}
--- schema violations become checkable. (2)~The standard alignment primitive
is at fault: matching facts by entity name hides most cross-model
disagreements. Anchoring each extracted figure to its unique source token
yields 1.7--2.8$\times$ as many comparable facts, and a quarter of anchored
values turn out to be unit-scale misreadings invisible to name matching.
(3)~Downstream, naive QA scoring says gold passages beat serialized graphs.
An \emph{evidence-conditional} evaluation --- separating memorization,
honest abstention, and grounded answers --- reverses that reading on two of
three models, attributes the residual gap to model-specific over-refusal,
and shows that attaching provenance pointers, which carry no content, raises
accuracy on the same two models. Every registered hypothesis carried its
disconfirming outcome in advance; two failed and are reported as findings.
Exploratory results are labelled as such and replicated under fresh
registration.
\end{abstract}

\section{Introduction}\label{sec:intro}

% From narrative/00-motivation + 01-problem, decision-B framing:
% industry -> domain -> ontology -> graph, then the diagnosis.

Enterprises that build knowledge graphs from documents with LLMs face a
design decision early: whether to guide extraction with a domain ontology.
The folk answer is yes --- a schema should make extraction more consistent,
more complete, and easier to query. Financial services is the domain where
this belief is most institutionalized: FIBO~\citep{fibo}, an OWL ontology of
financial concepts maintained by the EDM Council, exists precisely to
standardize how institutions name and relate the entities in their filings.

This paper measures that belief end to end on FinDER~\citep{finder}, a
financial question-answering dataset built from SEC filings, and reports
three findings that reorder it.

\textbf{First, ontology guidance does not improve extraction --- it makes bad
extraction detectable.} Across a five-level schema spectrum (from a
schema-free single-class floor to full FIBO with synonyms and subsumption),
every ontology-guided condition agrees \emph{less} across extractor models on
entity names than the schema-free floor, and the full-FIBO condition covers
\emph{fewer} of the figures that gold answers require ($-0.036$, 95\% CI
$[-0.067, -0.008]$, paired over 280 cases). What the ontology measurably adds
is a contract to violate: 529--1,038 schema violations per condition become
mechanically checkable, where the schema-free graph has almost nothing to
check.

\textbf{Second, the field's alignment primitive --- the entity name --- hides
the disagreements that matter.} Whether two extractors captured ``the same
fact'' is decided, in GraphRAG pipelines and their evaluations, by matching
names. Names are the least stable thing an LLM produces. Anchoring each
extracted figure to the unique numeric token it came from in the source
document multiplies comparable cross-model facts by 1.7--2.8$\times$ at the
same sample, and exposes that 26.7\% of anchored figures differ from their
source by a factor of $10^3$ or $10^6$ --- unit-scale misreadings
(``\$1,906,715'' extracted from a table headed ``in thousands'') that no
name-based comparison can see, because the two views also named the fact
differently. We call this \emph{provenance-anchored alignment} and replicate
its advantage on two extraction conditions at 280 cases each, pre-registered.

\textbf{Third, downstream evaluation inherits the same blindness.} Under
standard reference-based scoring, handing a model the source passages beats
handing it the serialized graph built from those passages --- the familiar
``RAG beats GraphRAG'' reading~\citep{ragvsgraphrag}. But public filings are
in every model's training data, so a correct answer is not evidence the
evidence was used. Conditioning each verdict on whether the served evidence
could support the answer separates grounded correctness from memorization
and honest abstention from failure. Under that decomposition the graph
conditions yield \emph{as many or more} genuinely grounded correct answers
than gold passages on two of three models, the text condition carries
2--4$\times$ the memorization, and the third model's graph deficit is
over-refusal --- it declines even when the graph contains the answer.

The through-line is one sentence: \emph{the value of structure is not a
better answer but a verifiable one, and verification needs provenance, not
names}. Every claim above was pre-registered with its disconfirming outcome
before the run that tested it (Section~\ref{sec:design}); two registered
hypotheses were disconfirmed and are reported as such.

\section{Related work}\label{sec:related}

% Ontology-guided extraction; GraphRAG evaluations; entity resolution;
% contamination-aware evaluation; LLM-as-judge practice; finance QA metrics.
\pending{related work — full pass after skeleton merge}

\section{Study design}\label{sec:design}

\subsection{Dataset and sampling}

FinDER pairs terse analyst questions with gold answers and the filing
passages that support them (5{,}703 cases, eight categories). Its own
annotations partition the corpus into three strata: \emph{simple} questions
(neither annotation), \emph{arithmetic} questions (an operation label such
as Division), and \emph{reasoning} questions (flagged as requiring
composition, no arithmetic). We run three samples that together cover the
grid: a category-balanced 280-case sample (35 per category, seed 42,
multi-reference preferred), a 140-case arithmetic-stratum sample
(type-proportional), and \pending{a 204-case census of the reasoning
stratum (an3)}. The balanced sample contains zero arithmetic cases --- the
multi-reference preference excluded them, since 881 of 883 arithmetic
questions carry one reference --- which is itself a finding about sample
design and the reason the supplement exists.

\subsection{The schema spectrum}

Five extraction conditions differ only in the schema handed to the
extractor: \textbf{A} a schema-free floor (one class, one relation);
\textbf{B} a hand-written 20-class schema; \textbf{C} 70 classes derived
from FIBO, corpus-scoped; \textbf{D} = C plus FIBO synonyms and
abbreviations; \textbf{E} = C plus the subsumption hierarchy. Everything
else --- documents, chunking, extractor models (DeepSeek-V3.1, gpt-oss-120b,
MiniMax-M2.7), prompts, property floor, graph store --- is held fixed, so
schema is the only moving part. Conditions A and C are re-run at 280 cases;
each case's graph is extracted independently per model, giving three
\emph{views} per case and condition.

\subsection{Pre-registration}

Each run is governed by a registration committed before it starts, stating
direction and the outcome that would disconfirm it; exploratory findings are
labelled and their registered replications named. The hypothesis ledger and
all registrations ship in the supplementary material. We report every
registered hypothesis, including the two that failed.

\section{Part 1: what the ontology changes in the graph}\label{sec:part1}

\subsection{Name agreement falls as the schema grows}

% SW-H4, S2-H2 — numbers final
Cross-model agreement on entity names is highest under the schema-free
floor and falls under every ontology condition (paired A$-$C difference in
per-case comparable-name rate: $+0.069$, 95\% CI $[+0.046, +0.091]$, 280
cases). The mechanism is structural: one declared class gives two extractors
nothing to disagree about; seventy classes are seventy naming choices.

\subsection{Content coverage does not rise}

% S2-H4 reversal — numbers final
The registered hypothesis that FIBO buys content at the price of names ---
supported at 16 cases --- \emph{reversed} at 280: condition C covers fewer
gold-answer figures than A ($-0.036$ $[-0.067, -0.008]$). The ontology has
no measured extraction benefit on this corpus.

\subsection{What survives: detectability}

Schema violations (missing required properties, undeclared labels, datatype
breaks) number 529--1,038 per ontology condition and are mechanically
checkable against the declared schema; the schema-free floor offers almost
nothing to check. The ontology's contribution relocates from generation
quality to \emph{verification surface} --- consistent with keeping heavy
ontology reasoning out of the extraction path and using it as an offline
validator.

\subsection{Why the synonym condition could not matter}

Only 30 of 5{,}703 questions (0.5\%) require vocabulary mapping at all, and
the abbreviations questions actually use (\texttt{rev}, \texttt{mgmt}) are
analyst shorthand, not FIBO designations. Condition D was measuring a need
the corpus does not have --- a caution for scoping ontology vocabulary work
by measured query demand rather than by ontology completeness.

\section{Provenance-anchored alignment}\label{sec:anchor}

% The paper's center. SC-H4, S2-H1 — numbers final.
Two views of the same case hold ``the same fact'' if we can align their
facts. The standard key is the entity name. We instead locate each extracted
figure at the unique numeric token it came from in the source passage
(searching all unit scales, so a mis-scaled extraction still anchors), and
align facts that anchor to the same token.

At 280 cases on the schema-free condition, name keys yield 992 comparable
fact pairs with 324 disagreements; anchor keys yield 1{,}656 pairs with 606
disagreements, of which \textbf{485 are invisible to any name key} --- the
views named the fact differently \emph{and} disagreed on its value. On the
FIBO condition the multiplier is larger (717 vs.\ 1{,}983 pairs; 246 vs.\
874 disagreements; 734 name-blind). Both replications were registered in
advance. The disagreements the anchor exposes are dominated by unit-scale
misreadings: 26.7\% of anchored figures match their source only after
rescaling by $10^3$ or $10^6$.

Two consequences follow. For pipeline builders: entity resolution keyed on
names silently discards most of the verification signal that running a
second extractor was supposed to buy. For evaluators: cross-model
verification has a hard ceiling regardless of key --- 87\% of facts are held
by exactly one view, a share that \emph{rises} with sample size (77.1\% at
16 cases) --- and agreement-gating trades four correct answers for every
wrong answer it blocks (precision $0.951 \to 1.000$, recall $0.951 \to
0.764$; the registered hypothesis that gating helps was disconfirmed on its
own criterion).

\section{Part 2: does the graph help answer?}\label{sec:part2}

\subsection{Design}

Each question is answered under five evidence conditions --- closed book,
gold passages, serialized schema-free graph, serialized FIBO graph, and the
FIBO graph with provenance pointers --- by the same three models at
temperature 0. The serializer is one code path; document-hierarchy nodes are
excluded because they carry the passages verbatim. Passages are the
\emph{ceiling} of retrieval (the gold evidence in full), and the graph
conditions serve the union of all three extractor views, the ceiling of what
extraction captured. Primary metric: scale-normalized overlap between the
gold answer's figures and the produced answer's (the finance-QA standard
family~\citep{finqa,tatqa}); a three-judge cross-judging panel (no judge
grades its own model's answers) covers what number matching cannot.

\subsection{The gate: contamination is real and stratified}

Closed-book performance spans $0.058$ (DeepSeek) to $0.306$ (gpt-oss) on the
balanced sample --- public filings are memorized to very different degrees
--- yet passages beat closed book for every model (paired differences
$+0.092$ to $+0.299$, all separated). On the arithmetic stratum the gate
blows open ($+0.41$ both models run): \emph{computed} figures cannot be
recalled, making arithmetic questions the contamination-resistant stratum
for evaluating retrieval on public corpora.

\subsection{Naive scoring: text wins or ties; format effects are model
properties}

Under naive overlap no graph condition beats passages anywhere; DeepSeek
loses with graphs on both ($+0.075$, $+0.062$ toward passages), gpt-oss on
the schema-free graph, MiniMax ties. The passages$-$graph gap itself differs
across models beyond resampling (DeepSeek vs.\ MiniMax: $+0.081$ $[+0.019,
+0.144]$) --- averaging models erases a real interaction, which is one
reason published RAG-vs-GraphRAG aggregates disagree.

\subsection{Evidence-conditional evaluation reverses the reading}

Scoring against the gold answer alone cannot distinguish a model that used
the evidence from one that remembered the filing. Because we can check
mechanically whether the served evidence contains the gold figures, every
verdict decomposes into: grounded-correct, utilization failure,
over-refusal, honest abstention, contaminated-correct, and hallucination.

\begin{table}[t]
\centering\small
\caption{Grounded-correct counts (evidence contains the answer and the model
answered it) on the balanced sample, 199 numeric-gold cases.}
\label{tab:ec}
\begin{tabular}{lcccc}
\toprule
 & passages & graph (schema-free) & graph (FIBO) & FIBO+pointers \\
\midrule
gpt-oss-120b & 107 & 111 & 105 & 112 \\
MiniMax-M2.7 & 124 & 138 & 127 & \textbf{146} \\
DeepSeek-V3.1 & 98 & 83 & 82 & 104 \\
\bottomrule
\end{tabular}
\end{table}

Three facts move (Table~\ref{tab:ec}). The graphs' grounded yield matches or
beats passages on two of three models. Passages carry 2--4$\times$ the
contamination (20--22 vs.\ 5--10 answers scored correct on figures absent
from the evidence). And DeepSeek's graph deficit is over-refusal --- 55
declines where the graph \emph{contains} the answer, vs.\ 31 with passages
--- a legibility failure of the serialization for that model, not a failure
of the graph's content. \pending{registered replication on the arithmetic
stratum: AR-EC1--3, an2}

\subsection{Provenance pointers raise accuracy without carrying content}

Attaching each anchored figure's source pointer (passage index and character
offset --- never the source text) was registered to change attribution, not
accuracy. The registration was disconfirmed in the interesting direction:
accuracy \emph{rose} on MiniMax ($+0.055$ $[+0.012, +0.097]$) and DeepSeek
($+0.050$ $[+0.022, +0.080]$). Since no answer content was attached, the
pointer acts as a trust signal --- it marks exactly the figures that
anchored to source --- and the pointer condition shows the highest evidence
grounding rate (73\%/70\% of answered figures present in the served
evidence). Mechanical citation checking lands at 42--58\% verified, so
pointer-following is far from solved; but provenance decoration is a
zero-content prompt intervention that buys accuracy.

\subsection{The arithmetic stratum: the original hypothesis meets its test}

\pending{AR-H2/H3/H4 + AR-EC1--3 verdicts, an2 — lands before the PDF
deadline}

\section{Discussion: from verdicts to design rules}\label{sec:discussion}

% The verdict -> cause -> prevention table, each row a measured count.
Owning the full chain --- source token, extracted node, served evidence,
answer, verdict --- lets each failure class be traced to its cause and
converted to a rule. Unit-scale misreads (26.7\%) argue for anchor-verifying
values at write time. Name fragmentation growing with schema size argues for
provenance keys in entity resolution and for deploying ontologies as
validators rather than extraction guides. The 87\% single-view ceiling caps
what multi-extractor redundancy can verify --- the budget belongs in
provenance. Model-specific over-refusal on serialized graphs makes
serialization a per-model design surface. And contamination that survives
inside evidence conditions makes evidence-conditional evaluation, not
closed-book gates alone, the honest default for public-corpus benchmarks.

\section{Limitations}\label{sec:limitations}

Single corpus and domain; three extractor models; passages are a retrieval
ceiling, not a retriever; the graph conditions test the graph as serialized
context, not as a queryable store (deferred, with the token-budget/cache
tension, to future work); ``evidence contains the answer'' is a numeric-token
approximation; per-type strata below $n=20$ are reported but underpowered.
\pending{judge panel results and inter-judge agreement}

\section{Conclusion}\label{sec:conclusion}

Ontology guidance, measured end to end, buys detectability rather than
extraction quality; the alignment primitive that decides what ``the same
fact'' means should be provenance, not names; and evaluation that cannot see
evidence cannot see what structure is for. The graph's contribution to
answering is not more correct answers --- it is answers that are grounded,
attributable, and honest about what the evidence cannot support.

\bibliographystyle{unsrtnat}
\bibliography{references}

\appendix
\section{Registrations and verdicts}
\pending{ledger table: every registered hypothesis, verdict, artifact}

\section{Scoring-metric sensitivity}
\pending{token-F1 table + the two tokenizer defects the smoke caught}

\end{document}
